\section{Google, OpenAI y Search: cómo se redistribuye el control de la puerta de entrada}

La sección anterior mostró que Google y OpenAI están en la mesa principal; esta sección analiza específicamente la competencia entre ambos por la puerta de entrada. El juicio más tenso del programa es este: el regreso triunfal de Google no significa que la crisis haya pasado, y ChatGPT tampoco es solo una herramienta de chat marginal. Al final, ambos podrían enfrentarse directamente en la búsqueda, los asistentes personales, los dispositivos móviles y los presupuestos publicitarios.

\subsection{Tras el ascenso de Gemini, ¿por qué OpenAI sigue siendo fuerte?}

Guangmi (广密) tiene una visión más bien optimista de OpenAI, por tres razones. Primera: OpenAI podría ser una de las pocas empresas de los últimos veinte años con una posibilidad real de desafiar a Google. Ni siquiera Bing, de Microsoft, cambió de fondo la puerta de entrada de la búsqueda; ChatGPT, en cambio, trasladó la mentalidad del usuario de «buscar respuestas» a «que la IA me ayude a entender, generar y ejecutar». Segunda: OpenAI sigue siendo uno de los mejores activos de IA: su densidad de talento, su sentido de producto, su marca y su escala de usuarios son muy sólidos. Tercera: OpenAI conserva una probabilidad alta de participar en la exploración del siguiente paradigma, lo cual es decisivo para su valuación a largo plazo.

Las ventajas de Google son igual de reales: activos como Gemini, la multimodalidad, Android, Search, YouTube, Workspace, Cloud, las TPU y Waymo forman una base muy amplia. El problema es que la fortaleza de Google no elimina automáticamente a OpenAI; este podría ser un mercado donde coexistan a largo plazo dos o incluso varios gigantes.

\subsection{Cómo reescribirá el Agent a Search}

Si, como se dijo antes, la crisis de Google no ha pasado realmente, la cuestión recae en Search: la búsqueda no es una función aislada, sino la puerta de entrada general a la publicidad, la distribución de páginas web, el tráfico hacia los comercios y las decisiones de los usuarios. Esta sección pone lado a lado el flujo de la búsqueda tradicional y el flujo de un Agent; solo así se ve por qué la redistribución del tráfico toca el núcleo de las utilidades de Google.

El flujo tradicional de la búsqueda es el siguiente: el usuario escribe palabras clave, el motor de búsqueda devuelve enlaces, resúmenes y anuncios, y el usuario hace clic, compara, compra o sigue buscando. El flujo de un Agent podría ser, en cambio: el usuario plantea un objetivo y el Agent busca, compara, invoca herramientas, negocia el precio, hace el pedido, paga, da seguimiento y, al final, devuelve el resultado. La diferencia entre ambos no es solo de interfaz, sino de posición en la cadena de valor.

\begin{figure}[H]
\centering
\includegraphics[width=0.96\textwidth]{search-agent-disruption.png}
\caption{Redistribución del tráfico entre Search y el Agent: de los enlaces de búsqueda a la tarea completada. Diagrama conceptual propio, elaborado a partir de la conversación de 00:31:20--00:35:08 y 00:50:57--00:52:53.}
\end{figure}

\begin{knowledgebox}{Lectura de la figura: lo que el Agent disputa es el poder de distribución}
En la figura, el objetivo del usuario entra primero al Agent, y no directamente al cuadro de búsqueda ni a una Super App. Después, el Agent decide si invoca la búsqueda, la venta de boletos, el transporte, los pagos o los sistemas empresariales. Quien controle esta puerta de entrada tendrá la oportunidad de redistribuir la publicidad, las transacciones, las comisiones y los datos de los usuarios.
\end{knowledgebox}

\begin{warningbox}{El Agent no eliminará la búsqueda de un momento a otro}
Para que el Agent reescriba de verdad la publicidad en buscadores, debe resolver la confiabilidad de las citas, la información en tiempo real, los pagos, los permisos, la integración de los comercios, la prevención de fraudes, la auditoría y la asignación de responsabilidades. A corto plazo es más probable una forma híbrida: la búsqueda tradicional incorpora IA, ChatGPT incorpora la búsqueda, y cada uno va erosionando el terreno del otro.
\end{warningbox}

\subsection{La mentalidad móvil de ChatGPT y la mentalidad de ecosistema de Gemini}

A continuación hay que pasar del mecanismo de la puerta de entrada a la mentalidad del usuario, porque modelos que responden igual de bien pueden terminar abriéndose en momentos, dispositivos y escenarios completamente distintos. Esta sección trata de la frecuencia de uso y del lugar que ocupa cada uno en la vida diaria, no solo de la capacidad del modelo.

En el programa se menciona una observación con un fuerte enfoque de producto: ChatGPT se parece más a un asistente para la vida diaria, que muchos usuarios usan los fines de semana, en el celular y en lo cotidiano; Gemini, en algunos escenarios, se parece más a una herramienta de productividad para el horario laboral, y cuenta con una ventaja natural de distribución gracias a Android y a los mercados emergentes. Esta diferencia implica que, además de la capacidad del modelo, la frecuencia de uso, la cercanía del escenario y la mentalidad del usuario importan igual.

Si ChatGPT logra convertirse en un asistente personal con 1,000--2,000 millones de usuarios activos diarios, su base comercial será muy estable; si Gemini logra integrar de verdad Android, Search, Workspace y la multimodalidad, Google también podría mantener su fortaleza en la era de la IA. Los dos caminos no se excluyen entre sí: compiten por cuál puerta de entrada abre primero el usuario.

\subsection{Resumen de la sección}

La crisis de Google pasó de «¿se le escapará la IA?» a «¿podrá conservar el poder de distribución en la puerta de entrada del Agent?». El reto de OpenAI pasó de «¿podrá construir buenos modelos?» a «¿podrá convertir ChatGPT en un asistente personal y una puerta de entrada de búsqueda y ejecución de uso frecuente, confiable y monetizable?». La conclusión de esta sección es que la guerra por la puerta de entrada apenas comienza, y los rankings de modelos son solo una pequeña parte de ella.
